\chapter{The laser pulse and tunnel ionization}
\label{ch:adk}

\section{The pulse}

The pulse is a linearly polarized Gaussian with a carrier wave,
\begin{equation}
  I(t) = I_0\,e^{-(t-t_0)^2/2\sigma^2},\qquad
  E(t) = E_0\,e^{-(t-t_0)^2/4\sigma^2}\cos\!\big(\omega(t-t_0)+\varphi\big),
\end{equation}
with $t_0$ = \opt{sfmd-pulse-peak-time}, $\sigma$ = \opt{sfmd-pulse-fwhm}$/2\sqrt{2\ln2}$, $\omega = 2\pi c/\lambda$
($\lambda$ = \opt{sfmd-pulse-wavelength}) and a carrier phase $\varphi$ drawn uniformly in $[-\pi/2,\pi/2]$ per run (and
broadcast, so that all ranks see the same field).
The peak intensity follows from the pulse energy $E_p$ = \opt{sfmd-pulse-energy} spread uniformly over a disc of diameter
$d$ = \opt{sfmd-pulse-focal-diameter}:
\begin{equation}
  I_0 = \frac{E_p}{\pi(d/2)^2\,\sigma\sqrt{2\pi}},\qquad E_0 = \sqrt{2I_0/(c\varepsilon_0)}.
\end{equation}
The sample (a few nm) is far smaller than the wavelength, so the field is uniform over it.
It has no direction in the model: ADK depends only on $|E(t)|$, and the cloud sees a cycle-averaged field
(Chapter~\ref{ch:laser}).
The field is evaluated within $\pm6\sigma$ of the peak and is zero outside; \file{pulse\_profile.txt} has $I(t)$ and
$E(t)$ every step.

\section{The ADK rate}

An atom in charge state $q$ is ionized to $q+1$ at the static-field tunnelling rate of Ammosov, Delone and
Krainov~\cite{adk1986}, in the form of Tong and Lin~\cite{tong2002} (atomic units):
\begin{gather}
  W = |C_{n^*l^*}|^2\,f(l,m)\,I_p\,\Big(\frac{2(2I_p)^{3/2}}{F}\Big)^{2n^*-|m|-1}
      \exp\!\Big(-\frac{2(2I_p)^{3/2}}{3F}\Big),\\
  |C_{n^*l^*}|^2 = \frac{1}{2\pi n^*}\Big(\frac{2e}{n^*}\Big)^{2n^*},\qquad
  f(l,m) = \frac{(2l+1)\,(l+|m|)!}{2^{|m|}\,|m|!\,(l-|m|)!},\qquad n^* = \frac{Z}{\sqrt{2I_p}},
\end{gather}
with $F = |E(t)|$ the instantaneous field, $I_p$ the ionization energy of the state, and $Z = q+1$ the charge left
behind.
The instantaneous (not cycle-averaged) rate is used, with the field of the current step, so ionization happens near the
field crests; at 2\,as there are $\sim$1400 steps per optical cycle at 850\,nm.
The constants of every state are computed once at start-up and printed to the \file{.log}.

\begin{table}[H]
\centering\small
\begin{tabular}{@{}l l l l@{}}
\toprule
element & ionization energies (eV), $q\to q+1$ & orbital ($n,l,|m|$) & states \\
\midrule
H & 13.598 & 1s & 1 \\
C & 11.260, 24.383, 47.887, 64.494 & 2p ($|m|$ 0, 1), then 2s & 4 \\
N & 14.534, 29.601, 47.448, 77.473, 97.890 & 2p ($|m|$ 0, 1, 1), then 2s & 5 \\
O & 13.618, 35.117, 54.936, 77.413, 113.899, 138.119 & 2p ($|m|$ 0, 1, 1, 1), then 2s & 6 \\
S & 10.360, 23.330, 34.830, 47.300, 72.600, 88.000 & 3p ($|m|$ 0, 1, 1, 1), then 3s & 6 \\
\bottomrule
\end{tabular}
\caption{The charge states ADK can reach: the valence shell only. The first electron comes from $m=0$, later ones from
$|m|=1$, a simplified Hund order.}
\label{tab:adk}
\end{table}

Only the valence shell is stripped: H to 1+, C to 4+, N to 5+, O and S to 6+.
Inner shells (K for C, N, O; L for S) are out of reach of optical fields at these intensities.
Other elements (F, Na, Mg, P, Cl, Ca, Fe, \dots) are recognized by their mass but have no ADK data and stay neutral.
Deuterium (mass 2) is hydrogen.

\section{Sampling}

Every MD step, for every atom, the rate $W$ of its current state gives the probability
\begin{equation}
  p = 1 - e^{-W\Delta t}
\end{equation}
of one ionization in the step, tested against a uniform random number from the module's own stream.
At most one electron per atom and step is released; at the rates reached here ($W\Delta t\ll1$ except at the field
crests at $10^{17}$\,W/cm$^2$, where states saturate within a few steps) that is no limitation.
Without the cloud every rank samples its own atoms with its own stream; with the cloud the master samples all atoms
from one stream, because each released electron's escape depends on the whole sample (Chapter~\ref{ch:pb}).
\file{transition\_log.txt} records every event (Chapter~\ref{ch:output}).

\section{What ADK leaves out}

\begin{itemize}
  \item Multiphoton ionization: ADK is the tunnelling limit, valid for a Keldysh parameter
    $\gamma = \omega\sqrt{2I_p}/F\lesssim1$; at 850\,nm and $10^{15}$\,W/cm$^2$, $\gamma\approx0.3$ for neutral C.
  \item Barrier suppression: above the over-the-barrier field ADK overestimates the rate; the states saturate anyway.
  \item The field of the neighbouring ions and of the cloud, which lowers or raises the barrier locally
    (enhanced ionization).
  \item Excited states, rescattering, and the electron's drift energy after the pulse: electrons are born at rest
    at their atom.
\end{itemize}
